\chapter{مبانی نظری}
\label{ch:theory}

\section{مقدمه}

در این فصل، مفاهیم بنیادی موردنیاز برای درک مسئله مجموعه احاطه‌گر، شبکه‌های عصبی گرافی و یادگیری تقویتی بیان می‌شود. ابتدا مفاهیم نظریه گراف مرور شده، سپس یادگیری تقویتی و در نهایت شبکه‌های عصبی گرافی معرفی می‌گردد.

\section{مبانی نظریه گراف}

نظریه گراف شاخۀ مهمی از ریاضیات گسسته است که برای مدلسازی روابط بین اشیاء به‌کار می‌رود. یک گراف ساده $G=(V,E)$ متشکل از مجموعه‌ای متناهی از رأس‌ها $V$ و مجموعه‌ای از یال‌ها $E \subseteq \{\{u,v\} \mid u,v \in V, u \neq v\}$ است \cite{bondy2008}.

\subsection{تعاریف پایه}

\begin{definition}[همسایگی]
مجموعه همسایگان رأس $v$ که با $N(v)$ نشان داده می‌شود، شامل تمام رأس‌های مجاور با $v$ است. همسایگی بسته $N[v]=N(v)\cup\{v\}$ تعریف می‌شود \cite{bondy2008}.
\end{definition}

\begin{definition}[درجه رأس]
درجه رأس $v$ برابر با $|N(v)|$ و با $deg(v)$ نمایش داده می‌شود \cite{bondy2008}.
\end{definition}

\subsection{مجموعه احاطه‌گر و مجموعه احاطه‌گر کامل}

\begin{definition}[مجموعه احاطه‌گر]
مجموعه $D \subseteq V$ یک مجموعه احاطه‌گر است اگر برای هر $v \notin D$، حداقل یک $u \in D$ وجود داشته باشد به‌طوری‌که $\{u,v\} \in E$؛ یعنی $\bigcup_{u\in D} N[u] = V$ \cite{haynes1998fundamentals}.
\end{definition}

\begin{definition}[مجموعه احاطه‌گر کامل]
مجموعه $T \subseteq V$ یک مجموعه احاطه‌گر کامل است اگر هر رأس $v \in V$ حداقل یک همسایه در $T$ داشته باشد؛ یعنی $\bigcup_{u\in T} N(u)=V$ \cite{haynes1998fundamentals}.
\end{definition}

تفاوت اساسی این است که در مجموعه احاطه‌گر کامل، رأس‌های عضو $T$ نیز باید توسط رأس دیگری از $T$ احاطه شوند؛ بنابراین هیچ رأس تنهای در $T$ نمی‌تواند نیاز همسایه‌اش را برطرف کند. این شرط، مسئله را برای کاربردهایی که نیاز به پوشش دوطرفه یا نظارت متقابل دارند، مناسب‌تر می‌سازد.

\section{یادگیری تقویتی}

یادگیری تقویتی\footnote{Reinforcement Learning (RL)} شاخه‌ای از یادگیری ماشین است که در آن یک عامل\footnote{Agent} از طریق تعامل با محیط\footnote{Environment} یاد می‌گیرد تا تصمیم‌گیری بهینه انجام دهد. در هر گام $t$، عامل حالت $s_t$ را مشاهده کرده، عمل $a_t$ را انتخاب می‌کند و پاداش $r_t$ دریافت می‌نماید \cite{sutton2018reinforcement}.

فرآیند تصمیم‌گیری مارکوف\footnote{Markov Decision Process (MDP)} معمولاً با پنج‌تایی $(S,A,P,R,\gamma)$ مدلسازی می‌شود. تابع ارزش $Q(s,a)$ مقدار مورد انتظار پاداش تجمعی آینده را پس از انتخاب عمل $a$ در حالت $s$ نشان می‌دهد.

\subsection{یادگیری Q عمیق دوگانه}

الگوریتم $DQN$ با استفاده از شبکه عصبی برای تقریب تابع $Q$ معرفی شد \cite{mnih2015human}. اما $DQN$ دچار بیش‌برآورد\footnote{Overestimation} مقادیر $Q$ می‌شود. برای رفع این مشکل، روش $Double\ DQN$\footnote{DDQN} پیشنهاد گردید که عمل انتخاب و ارزیابی ارزش را بین دو شبکه جداگانه تقسیم می‌کند \cite{vanhasselt2016ddqn}.

در $DDQN$، عمل با استفاده از شبکه آنلاین انتخاب و ارزش آن با شبکه هدف ارزیابی می‌شود:
$$
y_t = r_t + \gamma\, Q_{\theta^-}(s_{t+1}, \arg\max_a Q_\theta(s_{t+1},a))
$$
که $\theta$ پارامترهای شبکه آنلاین و $\theta^-$ پارامترهای شبکه هدف هستند.

\section{شبکه‌های عصبی گرافی}

شبکه‌های عصبی گرافی\footnote{Graph Neural Networks (GNNs)} خانواده‌ای از مدل‌های یادگیری عمیق هستند که مستقیماً روی داده‌های گرافی عمل می‌کنند. برخلاف شبکه‌های متعارف، GNNها قادرند ساختار توپولوژیکی گراف و اطلاعات رأس‌ها را به‌طور همزمان در نظر بگیرند \cite{scarselli2008graph, gilmer2017mpnn}.

\subsection{شبکه عصبی پیام‌رسان}

معماری رایج در GNNها، شبکه عصبی پیام‌رسان\footnote{Message Passing Neural Network (MPNN)} است. در هر لایه، هر رأس پیام‌هایی را از همسایگانش جمع‌آوری، آن‌ها را ترکیب و نمایش جدید خود را به‌روزرسانی می‌کند \cite{gilmer2017mpnn}.

فرآیند پیام‌رسانی به‌صورت زیر انجام می‌شود:
$$
m_v^{(k+1)} = \sum_{u \in N(v)} M_k(h_v^{(k)}, h_u^{(k)}, e_{vu})
$$
$$
h_v^{(k+1)} = U_k(h_v^{(k)}, m_v^{(k+1)})
$$
که $h_v^{(k)}$ نمایش رأس $v$ در لایه $k$، $M_k$ تابع پیام و $U_k$ تابع به‌روزرسانی است.

\section{خلاصه}

در این فصل مفاهیم پایه نظریه گراف، مسئله مجموعه احاطه‌گر و کامل، یادگیری تقویتی و شبکه‌های عصبی گرافی مرور شد. این مبانی، زیربنای روش پیشنهادی در فصل‌های بعدی را تشکیل می‌دهند.
